\documentclass[10pt,a4paper]{amsart}
\usepackage{amsmath,amssymb,amsthm,graphicx}
\usepackage{verbatim}
\newenvironment{reference}{\comment}{\endcomment}
%\newenvironment{reference}{}{}
\newenvironment{slogan}{\comment}{\endcomment}
\newenvironment{history}{\comment}{\endcomment}

% For commutative diagrams we use Xy-pic
\usepackage[all]{xy}

% We use 2cell for 2-commutative diagrams.
\xyoption{2cell}
% Theorem environments.
%
\theoremstyle{plain}
\newtheorem{theorem}[subsection]{Theorem}
\newtheorem{proposition}[subsection]{Proposition}
\newtheorem{lemma}[subsection]{Lemma}

\theoremstyle{definition}
\newtheorem{definition}[subsection]{Definition}
\newtheorem{example}[subsection]{Example}
\newtheorem{exercise}[subsection]{Exercise}
\newtheorem{situation}[subsection]{Situation}

\theoremstyle{remark}
\newtheorem{remark}[subsection]{Remark}
\newtheorem{remarks}[subsection]{Remarks}

\numberwithin{equation}{subsection}

% Macros
%
\def\lim{\mathop{\mathrm{lim}}\nolimits}
\def\colim{\mathop{\mathrm{colim}}\nolimits}
\def\Spec{\mathop{\mathrm{Spec}}}
\def\Hom{\mathop{\mathrm{Hom}}\nolimits}
\def\Ext{\mathop{\mathrm{Ext}}\nolimits}
\def\SheafHom{\mathop{\mathcal{H}\!\mathit{om}}\nolimits}
\def\SheafExt{\mathop{\mathcal{E}\!\mathit{xt}}\nolimits}
\def\Sch{\mathit{Sch}}
\def\Mor{\mathop{\mathrm{Mor}}\nolimits}
\def\Ob{\mathop{\mathrm{Ob}}\nolimits}
\def\Sh{\mathop{\mathit{Sh}}\nolimits}
\def\NL{\mathop{N\!L}\nolimits}
\def\CH{\mathop{\mathrm{CH}}\nolimits}
\def\proetale{{pro\text{-}\acute{e}tale}}
\def\etale{{\acute{e}tale}}
\def\QCoh{\mathit{QCoh}}
\def\Ker{\mathop{\mathrm{Ker}}}
\def\Im{\mathop{\mathrm{Im}}}
\def\Coker{\mathop{\mathrm{Coker}}}
\def\Coim{\mathop{\mathrm{Coim}}}

% Boxtimes
%
\DeclareMathSymbol{\boxtimes}{\mathbin}{AMSa}{"02}

%
% Macros for moduli stacks/spaces
%
\def\QCohstack{\mathcal{QC}\!\mathit{oh}}
\def\Cohstack{\mathcal{C}\!\mathit{oh}}
\def\Spacesstack{\mathcal{S}\!\mathit{paces}}
\def\Quotfunctor{\mathrm{Quot}}
\def\Hilbfunctor{\mathrm{Hilb}}
\def\Curvesstack{\mathcal{C}\!\mathit{urves}}
\def\Polarizedstack{\mathcal{P}\!\mathit{olarized}}
\def\Complexesstack{\mathcal{C}\!\mathit{omplexes}}
% \Pic is the operator that assigns to X its picard group, usage \Pic(X)
% \Picardstack_{X/B} denotes the Picard stack of X over B
% \Picardfunctor_{X/B} denotes the Picard functor of X over B
\def\Pic{\mathop{\mathrm{Pic}}\nolimits}
\def\Picardstack{\mathcal{P}\!\mathit{ic}}
\def\Picardfunctor{\mathrm{Pic}}
\def\Deformationcategory{\mathcal{D}\!\mathit{ef}}
\usepackage{cleveref}
\crefname{originalreference}{original reference}{original references}
\Crefname{originalreference}{Original reference}{Original references}
\makeatletter
\newlabel{lemma-pr2star}{{0FHC}{source}{}{sourcepointer-1}{}}
\newlabel{lemma-pr2star@cref}{{[originalreference][0][]0FHC}{[1][0][]source}{}{}{}}
\makeatother
\crefname{originalreference}{original reference}{original references}
\Crefname{originalreference}{Original reference}{Original references}
\makeatletter
\newlabel{lemma-trace-disjoint-union}{{0FHP}{source}{}{sourcepointer-2}{}}
\newlabel{lemma-trace-disjoint-union@cref}{{[originalreference][0][]0FHP}{[1][0][]source}{}{}{}}
\makeatother
\crefname{originalreference}{original reference}{original references}
\Crefname{originalreference}{Original reference}{Original references}
\makeatletter
\newlabel{lemma-unit}{{0FHF}{source}{}{sourcepointer-3}{}}
\newlabel{lemma-unit@cref}{{[originalreference][0][]0FHF}{[1][0][]source}{}{}{}}
\makeatother
\crefname{originalreference}{original reference}{original references}
\Crefname{originalreference}{Original reference}{Original references}
\makeatletter
\newlabel{lemma-weil-additive}{{0FHJ}{source}{}{sourcepointer-4}{}}
\newlabel{lemma-weil-additive@cref}{{[originalreference][0][]0FHJ}{[1][0][]source}{}{}{}}
\makeatother
\title{Cohomology of smooth projective schemes of dimension zero}
\author{The Stacks Project Authors}
\date{stacks-project:weil.tex:local-snapshot-20261012; GNU FDL 1.2 or later}
\begin{document}
\maketitle
\noindent\textit{Editorial note.} All delivered mathematical text is the author's own, copied from the pinned Stacks Project chapter file weil.tex; no statement or proof step was rewritten, reordered or removed.\par\medskip
\noindent\textit{Editorial note.} The counted proof is the author's complete argument: the Galois splitting computation of the cohomology of the zero dimensional scheme, the identification of the degree map with the trace, and the two statements about the fundamental class.\par\medskip
\noindent\textit{Editorial note.} The label lemma-dim-0 names the counted statement, and the proof environment immediately following it in the source is delivered verbatim as its proof.\par\medskip
\noindent\textit{Editorial note.} The counted statement is carried by the author's own lemma environment in the chapter source.\par\medskip
\noindent\textit{Editorial note.} Pre-open gate: python operations/agent-authoring/tools/precard\_gate.py --source weil.tex --label lemma-dim-0 --start 3431 --end 3502 --local-id 28128 returned RESULT=OK (L28128 free, occupancy FREE, public interval clear). Pre-assembly tag self-check: canonical tag resolved and not present in reference/public-items.jsonl.\par\medskip
\noindent\textit{Editorial note.} The statement and proof are copied from the pinned chapter file; the cited inputs are kept as external pointers with their canonical Stacks tags.\par\medskip
\noindent\textit{Allowed standard prerequisites (editorial).}\begin{itemize}
\item Additivity of the cohomology of disjoint unions and the structure of the cohomology of the base field. Same chapter, the additivity lemma for disjoint unions, the computation of the cohomology of the spectrum of a finite separable extension, and the nonvanishing of the cohomology of a nonempty scheme. Computes the cohomology of a product of the scheme with the spectrum of a finite Galois field by splitting the tensor product of the finite separable algebra into a product of fields.
\item Base change along a finite separable field extension. Axiom (B) on the behaviour of the theory under a finite separable base field extension. Identifies the cohomology of the product with the product of copies of the cohomology of the extension field.
\item The pushforward of the unit and the behaviour of the degree map over a product. Same chapter, the lemmas on the pushforward of the unit, the projection to the second factor, and the trace over a disjoint union. Identifies the degree map on the zero dimensional scheme with the trace of the corresponding separable algebra.
\end{itemize}\medskip
\ifdefined\hypertarget\hypertarget{sourcepointer-1}{}\fi
\noindent\textit{Source reference.} Original reference 0FHC: sources/stacks-weil/pinned-original-source.tex, source lines 2699--2699. Prerequisite pointer invoked at proof line(s) 3485; it names a standard fact of the same chapter that the counted argument uses, and every step of the counted proof built on it is retained verbatim in the delivered body.\par\medskip
\ifdefined\hypertarget\hypertarget{sourcepointer-2}{}\fi
\noindent\textit{Source reference.} Original reference 0FHP: sources/stacks-weil/pinned-original-source.tex, source lines 3408--3408. Prerequisite pointer invoked at proof line(s) 3481; it names a standard fact of the same chapter that the counted argument uses, and every step of the counted proof built on it is retained verbatim in the delivered body.\par\medskip
\ifdefined\hypertarget\hypertarget{sourcepointer-3}{}\fi
\noindent\textit{Source reference.} Original reference 0FHF: sources/stacks-weil/pinned-original-source.tex, source lines 2800--2800. Prerequisite pointer invoked at proof line(s) 3461; it names a standard fact of the same chapter that the counted argument uses, and every step of the counted proof built on it is retained verbatim in the delivered body.\par\medskip
\ifdefined\hypertarget\hypertarget{sourcepointer-4}{}\fi
\noindent\textit{Source reference.} Original reference 0FHJ: sources/stacks-weil/pinned-original-source.tex, source lines 2932--2932. Prerequisite pointer invoked at proof line(s) 3456; it names a standard fact of the same chapter that the counted argument uses, and every step of the counted proof built on it is retained verbatim in the delivered body.\par\medskip
\begin{lemma}
\label{lemma-dim-0}
Assume given (D0), (D1), (D2), and (D3) satisfying (A), (B), and (C).
Let $X$ be a smooth projective scheme of dimension zero over $k$.
Then
\begin{enumerate}
\item $H^i(X) = 0$ for $i \not = 0$,
\item $H^0(X)$ is a finite separable algebra over $F$,
\item $\dim_F H^0(X) = \deg(X \to \Spec(F))$,
\item $\int_X : H^0(X) \to F$ is the trace map,
\item $\gamma([X]) = 1$, and
\item $\int_X \gamma([X]) = \deg(X \to \Spec(k))$.
\end{enumerate}
\end{lemma}

\begin{proof}
We can write $X = \Spec(k')$ where $k'$ is a finite separable
algebra over $k$. Observe that $\deg(X \to \Spec(k)) = [k' : k]$.
Choose a finite Galois extension $k''/k$ containing each of the
factors of $k'$. (Recall that a finite separable $k$-algebra is
a product of finite separable field extension of $k$.)
Set $\Sigma = \Hom_k(k', k'')$. Then we get
$$
k' \otimes_k k'' = \prod\nolimits_{\sigma \in \Sigma} k''
$$
Setting $Y = \Spec(k'')$ axioms (B)(a) and Lemma \ref{lemma-weil-additive} give
$$
H^*(X) \otimes_F H^*(Y) =
\prod\nolimits_{\sigma \in \Sigma} H^*(Y)
$$
as graded commutative $F$-algebras. By Lemma \ref{lemma-unit} the
$F$-algebra $H^*(Y)$ is nonzero. Comparing dimensions on either side
of the displayed equation we conclude that $H^*(X)$ sits only in degree $0$
and $\dim_F H^0(X) = [k' : k]$. Applying this to $Y$ we get
$H^*(Y) = H^0(Y)$. Since
$$
H^0(X) \otimes_F H^0(Y) = H^0(Y) \times \ldots \times H^0(Y)
$$
as $F$-algebras, it follows that $H^0(X)$ is a separable $F$-algebra
because we may check this after the faithfully flat base change
$F \to H^0(Y)$.

\medskip\noindent
The displayed isomorphism above is given by the map
$$
H^0(X) \otimes_F H^0(Y) \longrightarrow
\prod\nolimits_{\sigma \in \Sigma} H^0(Y),\quad
a \otimes b \longmapsto \prod\nolimits_\sigma \Spec(\sigma)^*a \cup b
$$
Via this isomorphism we have $\int_{X \times Y} = \sum_\sigma \int_Y$ by
Lemma \ref{lemma-trace-disjoint-union}. Thus
$$
\int_X a = \text{pr}_{1, *}(a \otimes 1) = \sum \Spec(\sigma)^*a
$$
in $H^0(Y)$; the first equality by Lemma \ref{lemma-pr2star}
and the second by the observation we just made. Choose an
algebraic closure $\overline{F}$ and
a $F$-algebra map $\tau : H^0(Y) \to \overline{F}$.
The isomorphism above base changes to the isomorphism
$$
H^0(X) \otimes_F \overline{F} \longrightarrow
\prod\nolimits_{\sigma \in \Sigma} \overline{F},\quad
a \otimes b \longmapsto \prod\nolimits_\sigma \tau(\Spec(\sigma)^*a) b
$$
It follows that $a \mapsto \tau(\Spec(\sigma)^*a)$ is a full set
of embeddings of $H^0(X)$ into $\overline{F}$. Applying $\tau$
to the formula for $\int_X a$ obtained above we conclude
that $\int_X$ is the trace map.
By Lemma \ref{lemma-unit} we have $\gamma([X]) = 1$.
Finally, we have $\int_X \gamma([X]) = \deg(X \to \Spec(k))$
because $\gamma([X]) = 1$ and the trace of $1$ is equal to $[k' : k]$
\end{proof}

\end{document}
